\documentclass[a4paper,12pt]{article}

\usepackage[croatian]{babel}

\usepackage[a4paper,top=2.5cm,bottom=2.5cm,left=2.5cm,right=2.5cm,headheight=1.25cm,headsep=1.25cm,footskip=1.25cm]{geometry}
\usepackage[croatian,refpage]{nomencl}
\makenomenclature
\usepackage{fancyhdr}
\usepackage{tocloft}
\usepackage{lastpage}
\usepackage{amsmath}  
\usepackage{bm} 
\usepackage{tikz}
\usetikzlibrary{patterns,decorations.pathmorphing}
\usetikzlibrary{arrows.meta, positioning, shapes, decorations.markings}
\usepackage{tikz-cd}
\usepackage[labelfont=bf,textfont=bf]{caption}
\usepackage{amsfonts}
\usepackage[colorlinks=true, allcolors=black]{hyperref}
\usepackage{mathrsfs}
\usepackage{listings}
\usepackage{xcolor}
\usepackage{mathtools}
\usepackage{float}
\usepackage{caption}
\usepackage{csquotes}
\usepackage{longtable}
\usepackage{tikz-3dplot}
\usepackage{pgfplots}
\pgfplotsset{width=10cm,compat=1.18}
\usepackage{graphicx}
\usepackage{subcaption}
\usepackage{amsthm}

\usepackage{titlesec}

\titleformat{\section} % format \section
  {\bfseries\fontsize{18}{22}\selectfont} % bold + 18pt + 22pt leading
  {\thesection.} % section number with dot
  {1em} % space between number and title
  {} % before title code

\titleformat{\subsection} % optional: format \subsection
  {\bfseries\fontsize{16}{20}\selectfont}
  {\thesubsection.}
  {1em}
  {}

\titleformat{\subsubsection} % optional
  {\bfseries\fontsize{14}{18}\selectfont}
  {\thesubsubsection.}
  {1em}
  {}

\usepackage{etoolbox}

% Patch \tableofcontents to change title font
\patchcmd{\tableofcontents}
  {\section*{\contentsname}} % original
  {{\centering\bfseries\fontsize{20}{24}\selectfont \contentsname\par}} % new font
  {}{} % success/failure messages

% Theorem environments
\newtheorem{theorem}{Teorem}
\newtheorem{definition}[theorem]{Definicija}
\newtheorem{corollary}[theorem]{Korolar}
\newtheorem{proposition}[theorem]{Propozicija}
\newtheorem{remark}[theorem]{Napomena}

% Listings (Matlab)
\lstset{
    language=Matlab,
    basicstyle=\ttfamily\small,
    keywordstyle=\color{blue},
    commentstyle=\color{green!50!black},
    stringstyle=\color{orange},
    numbers=left,
    numberstyle=\tiny\color{gray},
    stepnumber=1,
    numbersep=5pt,
    showstringspaces=false,
    breaklines=true,
    frame=single,
    framerule=1pt,
    backgroundcolor=\color{gray!10}
}

% Dot after section numbers
\titlelabel{\thetitle.\quad}
\renewcommand{\cftsecpresnum}{}
\renewcommand{\cftsecaftersnum}{.}
\renewcommand{\cftsubsecpresnum}{}
\renewcommand{\cftsubsecaftersnum}{.}

% Header/footer style (optional, you can customize)
\usepackage{fancyhdr}

% Define custom style
\fancypagestyle{numbered}{
    \fancyhf{}
    \fancyhead[L]{\fontsize{12}{14}\selectfont Fakultet strojarstva i brodogradnje}
    \fancyhead[R]{\fontsize{12}{14}\selectfont Računalna dinamika fluida}
    \fancyfoot[L]{\fontsize{12}{14}\selectfont Pero Perić}
    \fancyfoot[C]{}
    \fancyfoot[R]{\thepage}
    \renewcommand{\headrulewidth}{0.4pt}
    \renewcommand{\footrulewidth}{0.4pt}
}

% Apply to all normal pages
\pagestyle{numbered}

% Make plain style (used by TOC, LOF, LOT, chapter* pages) the same as numbered
\fancypagestyle{plain}{
    \fancyhf{}
    \fancyhead[L]{\fontsize{12}{14}\selectfont Fakultet strojarstva i brodogradnje}
    \fancyhead[R]{\fontsize{12}{14}\selectfont Računalna dinamika fluida}
    \fancyfoot[L]{\fontsize{12}{14}\selectfont Pero Perić}
    \fancyfoot[C]{}
    \fancyfoot[R]{\thepage}
    \renewcommand{\headrulewidth}{0.4pt}
    \renewcommand{\footrulewidth}{0.4pt}
}



\begin{document}

%%%%%%%%%%%%%%%%%%%%%%%%%%%%%%%%%%%%%%%%%%%%%%%%%%%%%%%%%%%%%
% FRONT PAGE (same layout as your picture)
%%%%%%%%%%%%%%%%%%%%%%%%%%%%%%%%%%%%%%%%%%%%%%%%%%%%%%%%%%%%%

\begin{titlepage}
\begin{center}

{\fontsize{20}{28}\selectfont \text{FAKULTET STROJARSTVA I BRODOGRADNJE}}\\

\vspace*{0.5cm}
{\fontsize{18}{22}\selectfont SVEUČILIŠTE U ZAGREBU}

\vspace*{5cm}

{\fontsize{22}{24}\selectfont Seminarski rad iz kolegija}\\
\vspace*{1cm}
{\fontsize{24}{24}\selectfont \textbf{\textit{Računalna dinamika fluida}}}

\vspace{1cm}

{\fontsize{18}{18}\selectfont prof. dr. sc. Spužva Bob}\\
{\fontsize{18}{18}\selectfont mag. ing. Patrik Zvijezda}

\vfill

\begin{flushleft}
{\fontsize{14}{18}\selectfont Ime i prezime: \textit{Pero Perić}}\\
{\fontsize{12}{18}\selectfont Matični broj: 0000000000}
\end{flushleft}


\vspace{1.5cm}

\begin{flushright}
{\fontsize{14}{18}\selectfont Zagreb, 2026./2027.}
\end{flushright}



\end{center}
\end{titlepage}

%%%%%%%%%%%%%%%%%%%%%%%%%%%%%%%%%%%%%%%%%%%%%%%%%%%%%%%%%%%%%
% FRONT MATTER
%%%%%%%%%%%%%%%%%%%%%%%%%%%%%%%%%%%%%%%%%%%%%%%%%%%%%%%%%%%%%

\newpage
\pagenumbering{Roman}

\tableofcontents

\newpage
\listoffigures
\addcontentsline{toc}{section}{Popis slika}

\newpage
\listoftables
\addcontentsline{toc}{section}{Popis tablica}

\newpage
\section*{Popis oznaka}
\addcontentsline{toc}{section}{Popis oznaka}
% \printnomenclature % uncomment if you actually use nomenclatures

%%%%%%%%%%%%%%%%%%%%%%%%%%%%%%%%%%%%%%%%%%%%%%%%%%%%%%%%%%%%%
% MAIN MATTER
%%%%%%%%%%%%%%%%%%%%%%%%%%%%%%%%%%%%%%%%%%%%%%%%%%%%%%%%%%%%%
\newpage
\pagenumbering{arabic}

\section{Zadatak}
Navier-Stokes jednadžbe

\begin{equation}
\begin{cases}
\rho \left( \frac{\partial u}{\partial t} + u \frac{\partial u}{\partial x} + v \frac{\partial u}{\partial y} + w \frac{\partial u}{\partial z} \right) = - \frac{\partial p}{\partial x} + \mu \left( \frac{\partial^2 u}{\partial x^2} + \frac{\partial^2 u}{\partial y^2} + \frac{\partial^2 u}{\partial z^2} \right) + \rho f_x, \\[1em]
\rho \left( \frac{\partial v}{\partial t} + u \frac{\partial v}{\partial x} + v \frac{\partial v}{\partial y} + w \frac{\partial v}{\partial z} \right) = - \frac{\partial p}{\partial y} + \mu \left( \frac{\partial^2 v}{\partial x^2} + \frac{\partial^2 v}{\partial y^2} + \frac{\partial^2 v}{\partial z^2} \right) + \rho f_y, \\[1em]
\rho \left( \frac{\partial w}{\partial t} + u \frac{\partial w}{\partial x} + v \frac{\partial w}{\partial y} + w \frac{\partial w}{\partial z} \right) = - \frac{\partial p}{\partial z} + \mu \left( \frac{\partial^2 w}{\partial x^2} + \frac{\partial^2 w}{\partial y^2} + \frac{\partial^2 w}{\partial z^2} \right) + \rho f_z, \\[0.5em]
\frac{\partial u}{\partial x} + \frac{\partial v}{\partial y} + \frac{\partial w}{\partial z} = 0
\end{cases}
\end{equation}


\section{Ručni proračun}
Detaljno opisati ručni proračun sa svim jednadžbama i slikama.

\begin{equation}
\begin{cases}
\rho \left( \frac{\partial \mathbf{u}}{\partial t} + (\mathbf{u} \cdot \nabla) \mathbf{u} \right) = - \nabla p + \mu \nabla^2 \mathbf{u} + \rho \mathbf{f}, \\[0.5em]
\nabla \cdot \mathbf{u} = 0
\end{cases}
\end{equation}


\section{Proračun}
Detaljno opisati postupak rješavanja zadatka.\cite{lanczos1986}


\section{Usporedba rezultata}
Usporediti.\\

Einsteinove jednadžbe polja. 

\begin{equation}
R_{\mu\nu} - \frac{1}{2} g_{\mu\nu} R + g_{\mu\nu} \Lambda = \frac{8 \pi G}{c^4} T_{\mu\nu}, \quad \mu,\nu = 0,1,2,3
\end{equation}

\noindent gdje je Ricci tenzor definiran kao:

\begin{equation}
R_{\mu\nu} = \partial_\lambda \Gamma^\lambda_{\mu\nu} - \partial_\nu \Gamma^\lambda_{\mu\lambda} + \Gamma^\lambda_{\mu\nu} \Gamma^\sigma_{\lambda\sigma} - \Gamma^\sigma_{\mu\lambda} \Gamma^\lambda_{\nu\sigma}
\end{equation}

\noindent a skalarski Ricci tenzor je:

\begin{equation}
R = g^{\mu\nu} R_{\mu\nu}
\end{equation}

\noindent Christoffelovi simboli \(\Gamma^\lambda_{\mu\nu}\) su definirani preko metrike \(g_{\mu\nu}\) kao:

\begin{equation}
\Gamma^\lambda_{\mu\nu} = \frac{1}{2} g^{\lambda\sigma} \left( \partial_\mu g_{\nu\sigma} + \partial_\nu g_{\mu\sigma} - \partial_\sigma g_{\mu\nu} \right)
\end{equation}



\section{Zaključak}
Kratki sadržaj i zaključci analize.

\section{Dodatak}
Sve snimke zaslona, rezultati.

%%%%%%%%%%%%%%%%%%%%%%%%%%%%%%%%%%%%%%%%%%%%%%%%%%%%%%%%%%%%%
% LITERATURA
%%%%%%%%%%%%%%%%%%%%%%%%%%%%%%%%%%%%%%%%%%%%%%%%%%%%%%%%%%%%%
\newpage
\addcontentsline{toc}{section}{Literatura}
\bibliographystyle{unsrt}
\bibliography{literatura_RDF}

\end{document}
